% Source: https://artofproblemsolving.com/community/c3h2483605_nt_problems
% Stable identity: aops-topic-2483605-q08
Since $103$ is odd, $\left(\sqrt3-5\right)^{103}=-\left(5-\sqrt3\right)^{103}$, so the number is
\[N=\left(5+\sqrt3\right)^{103}+\left(5-\sqrt3\right)^{103}=2\sum_{j=0}^{51}\dbinom{103}{2j}5^{103-2j}\,3^{j},\]
because in the binomial expansions the terms with odd powers of $\sqrt3$ cancel, while $\left(\sqrt3\right)^{2j}=3^j$. So $N$ is an integer.

Modulo $9$, every term with $j\ge2$ contains the factor $3^j$, a multiple of $9$, so it vanishes. The term $j=1$ is $\dbinom{103}{2}\cdot5^{101}\cdot3$, and $\dbinom{103}{2}=5253=3\cdot1751$, so it is also a multiple of $9$. Hence
\[N\equiv2\cdot5^{103}\pmod9.\]
The powers of $5$ modulo $9$ repeat with period $6$ ($5,7,8,4,2,1$), and $103\equiv1\pmod6$, so $5^{103}\equiv5$ and $N\equiv10\equiv1\pmod9$. The remainder is $\boxed{1}$.
